\paragraph{Proof and binary.} The theorems concern the translation of the MIR of
rustc \fact{rustc-mir-nightly}; rule-core's CI also builds the evaluator with
rustc \fact{rustc-rule-core-stable}, Perch builds with rustc \fact{rustc-perch},
and nothing checks the Wasm against the translation. Allocation failure and stack
exhaustion, which the model omits, trap. The specification is trusted, including
\code{Doc.allows}, written from Perch's earlier OpenZeppelin-based account
(\fact{doc-allows-source}) and not re-derived from the pinned gate.

\paragraph{Shell and host.} Tested, not proved: the parser and compiler,
stored-value decoding, rule-ID resolution, context conversion, the excluded set's
assembly, record persistence, the credential index, the generation update, the
composition of evaluation with authentication, and the recovery controller. The
integration tests run them as natively registered contracts. The host models
omit TTL, archival, fees, footprints and pre-host validation.

\paragraph{Grantees (A1) and submitters (A2).} A key named by several live rules
has the union of their authority; caps are per rule, not per grant. Allowances
used by \code{transfer\_from}, and funds held by other contracts, move without
the account's authorization, and a token rule that does not restrict functions
also allows \code{approve}. An owner-authorized \code{execute} runs as the
account, so no rule or cap sees the callee's \code{require\_auth}; grants cannot
reach it, but an administrative rule that allows it has unlimited authority over
its targets. A grantee can fill its window with \fact{spends-per-window} small
spends, and replacing a capped rule or its cap starts an empty window. Delegated
signers do not bind rule IDs, a rule edited in place decides under its new
contents, a refused entry can be resubmitted until it expires, and grant installs
are unordered.

\paragraph{Owner-key theft and recovery (A3).} Under Loss a thief holding the
owner key can cancel every recovery; Protected without an enrolled baseline
protects no better. The owner may schedule an upgrade outside an attempt and
execute it after \fact{upgrade-delay-ledgers} ledgers
(\fact{upgrade-delay-days}); new code inherits the account's storage, including
the revoked set and generation, and the proofs do not cover it. A completion
marker is consumed by the target's canonical identity, not its exact bytes
(issue \fact{completion-marker-issue}).

\paragraph{Trusted components.} The account trusts its verifiers and the recovery
controller its document names; the controller gets no registry provenance check,
so wallets should pin its content address. Once pinned code is deployed, no
registry action changes it, but a verified name means only that the registry
deployed that code at its own hash: the manager (A5) decides which hashes it
deploys, its admin can replace the registry's code, and one key holds both roles
on testnet and mainnet. WebAuthn origin is not checked, and the wallet must show
what a signed payload commits to.

\paragraph{Resources, coverage and audit.} Bounded records and arithmetic can
refuse permitted actions. Documents are capped at \fact{doc-max-signers}
signers, \fact{doc-max-rules} rules, \fact{doc-max-canonical-bytes} canonical
bytes and \fact{doc-max-rule-name-bytes}-byte rule names, so a key thief cannot
make recovery too costly to complete; the worst case at those caps uses
\fact{worst-case-completion-share} of the instruction limit, and disk-read bytes
are not measured. \code{renew} does not extend installed-grant records. No single
theorem covers the whole authorization path, including mixed administrative and
grant batches (Proposition~\ref{prop:end-to-end} is argued), and no external
audit has been completed.
